\section{Introduction}

Sharing and fine-tuning pretrained models has become a default practice in modern machine learning. Public repositories such as Hugging Face host hundreds of thousands of models, and practitioners routinely download these weights, attach a fresh classification head, and fine-tune on their own data. The datasets used for fine-tuning frequently contain sensitive information, such as medical records, user conversations, or internal communications. It is therefore standard practice to analyze the privacy properties of the resulting model and, increasingly, to adopt differentially private training procedures such as DP-SGD~\citep{abadi2016deep}. 

Two independent classes of risk have been studied in the fine-tuning pipeline. The first is a \textbf{supply-chain} risk that stems from the source of the pretrained weights: a malicious provider can tamper with a model's weights so that it misbehaves on specific inputs, compromising the model's \emph{integrity}~\citep{gu2017badnets, liu2018trojaning, hong2022handcrafted}. The second is a \textbf{memorization} risk that stems from the training procedure itself: neural networks trained (or fine-tuned) on sensitive data encode non-trivial information about individual examples, enabling membership inference~\citep{shokri2017membership} and data extraction~\citep{carlini2021extracting}. DP-SGD is the canonical defense against the second risk, bounding the per-example influence of the final weights, and is typically deployed under the implicit assumption that the initialization is benign. \\

\paragraph{Privacy backdoors.} \citet{feng2024privacy} introduce \emph{privacy backdoors}, which bridge the integrity and privacy threats described above. They show that a malicious provider can tamper with a pretrained model's weights so that, during ordinary fine-tuning on the victim's private dataset, individual training examples are imprinted into the weights in a form that the attacker can later recover. The construction relies on a novel primitive, which the authors term a \emph{data trap}: a small number of units are configured so that the gradient update induced by a specific training example is written with high fidelity into an attacker-chosen set of coordinates, and so that no subsequent update overwrites the signal. An adversary with access only to the final fine-tuned model can then read off the captured example directly from those coordinates.

\paragraph{Implications for DP-SGD.} The most consequential aspect of the attack is arguably not the reconstruction of the individual examples but its bearing on \emph{differential privacy}. The privacy analysis of DP-SGD rests on a worst-case condition: the per-example gradient differs maximally, on some coordinate, between the neighboring datasets~\citep{abadi2016deep}. The induced per-step bound is tight for the Gaussian mechanism; its composition across training steps, however, is widely believed to be pessimistic in practice, as the corresponding adversary is assumed to observe intermediate noisy gradient~\citep{nasr2021adversary}. In contrast, a realistic \emph{end-to-end} adversary, with access only to the deployed fine-tuned model, is expected to extract substantially less information than this bound permits; an intuition that has motivated the use of comparatively loose privacy budgets in the production deployments (e.g., $\varepsilon \geq 8$, or $\varepsilon \approx 9$ in \citet{ramaswamy2020training}).

Privacy backdoors invalidate this intuition. The data-trap primitive can be specialized so that, at every DP-SGD step, the gradient of the target example concentrates entirely on a small, attacker-known index set, with $\ell_2$-norm equal to the clipping threshold, while the gradients of all other examples vanish on that set. This constitutes precisely the per-step worst case assumed by the analysis. \citet{feng2024privacy} shows analytically that, under this construction, an end-to-end adversary observing only the final weights can derive a lower bound on the realized $\varepsilon$ that closely matches the provable upper bound. The discrepancy between the theoretical and empirical privacy of end-to-end DP-SGD is thus not attributable to structural slack in the analysis; it is a consequence of  benign initializations and disappears under an adversarial one.

\paragraph{Scope.} This report provides a self-contained exposition of \citet{feng2024privacy}, organized around two questions.
\begin{enumerate}[label=(\roman*), leftmargin=*, itemsep=2pt, topsep=2pt]
  \item \textbf{Construction.} How can an adversary engineer a pretrained initialization whose weights capture individual fine-tuning examples and preserve the captured signal across many SGD steps with gradient clipping and Gaussian noise?
  \item \textbf{Tightness.} Why does such an adversarial initialization render the end-to-end privacy leakage of DP-SGD nearly tight against its provable upper bound, and what is the resulting closed-form lower bound on the realized $\eps$?
\end{enumerate}

